%% ===========================================================================
\FloatBarrier
\section{Parameter identification and limitations}
\label{sec:calibration}
%% ===========================================================================

Four parameters are supplied by no constitutive table: the window size $a$,
the fractional order $\alpha$, the relaxation time $\tau$ and the flexoelectric
coefficient $\mu$. This section states how each is obtained, over which ranges
the results hold, and what is not verified.

\FloatBarrier
\subsection{Measurement routes for the four parameters}
\label{sec:calib_a_mu}

The four are not of the same kind. The geometry of the interaction, $M_{mn}$
and the orientation $\theta$ of its principal axes, comes from size-effect and
texture measurements; the material response, $\alpha$ and $\tau$, from
dielectric loss; $\mu$ from a bending or a direct charge measurement. The
protocol is ordered so that each step fixes one quantity and the later steps
check the earlier ones, as in the characteristic-length identification of the
one-scale theory.

Table~\ref{tab:protocol} of \ref{app:data} lists the seven steps in the order
in which they fix the quantities. Steps 1, 3 and 4 determine $M_{22}$ from three
independent observables, so that agreement among them checks the second-moment
description instead of repeating one measurement; step~6 turns the calibrated
tensor into a prediction of Figs.~\ref{fig:orientation} and~\ref{fig:design};
step~7 is the test of the description itself and is not attempted here.


\paragraph{The flexoelectric coefficient} The direct route is a four-point
bending test on a poled bar, electrodes on the two faces: for the geometry of
Fig.~\ref{fig:geometry}(a) the open-circuit voltage at a known curvature
$\kappa_{c}$ is
\begin{equation}
  V = \frac{\mu\,\kappa_{c}\,h}{\kappa_{yy}},
  \label{eq:muV}
\end{equation}
with no discretization involved, $\kappa_{yy}=4.9583\times10^{-10}$\,\si{\farad
\per\metre} the absolute permittivity of \eqref{eq:D}. The inverse route
measures the charge released by a device whose curvature is known from its
deflection. Both fix the magnitude, linear in $\mu$; the sign comes from the
crystal symmetry. The value used here, $\mu=10^{-6}$\,\si{\coulomb\per\metre},
is nominal: every bending coupling of Secs.~\ref{sec:thickness}
and~\ref{sec:transient} scales with $\mu^{2}$ and is a shape, not a magnitude,
until a measured coefficient is substituted.

\FloatBarrier
\subsection{Identification of the fractional order and of the relaxation time}
\label{sec:calib_alpha_tau}

The pair $(\alpha,\tau)$ comes from a dielectric loss measurement. The
fractional Kelvin--Voigt element \eqref{eq:kvfreq} has
$\tan\delta=(\omega\tau)^{\alpha}\sin(\pi\alpha/2)/
[1+(\omega\tau)^{\alpha}\cos(\pi\alpha/2)]$, which increases monotonically to
$\tan(\pi\alpha/2)$ and has no maximum, and both $|K^{*}|/K$ and $\tan\delta$
are monotone in $\omega\tau$; so no single frequency carries $\tau$, and the
pair $(|K^{*}|/K,\tan\delta)$ over two decades is fitted with $(\alpha,\tau)$.
At $\omega\tau=1$ the loss is $\tan(\pi\alpha/4)$, which gives $\alpha$ from one
decade once $\tau$ is known.
The ring-down gives the same information in the time domain, but not from one
number: for small loss $\zeta=(\omega\tau)^{\alpha}\sin(\pi\alpha/2)/2$, so
$\delta\approx\pi(\omega\tau)^{\alpha}\sin(\pi\alpha/2)$ -- at $\alpha=1$,
$\delta=\pi\omega\tau$, which reproduces \eqref{eq:zeta} to the $1.2\,\%$ of
the small-loss approximation. A single decrement fixes only that product, so
the pair needs two frequencies or one decrement at a known $\omega\tau$.

Two cautions follow from Sec.~\ref{sec:memorystudy}: the calibration must say
which matrix the memory multiplies, since one measured decrement gives two
different $(\alpha,\tau)$ pairs under $K_{\rm eff}$ and $K_{uu}$ (their damping
ratios differ by the factor $0.450$ at $\alpha=1$); and the fit is valid only
where the measurement is -- a broadened single-power-law loss. Two well
separated loss peaks need two elements.

\FloatBarrier
\subsection{Validity window}
\label{sec:limit_window}

Three inequalities bound the results, and each is a modeling assumption rather
than a numerical tolerance.

\begin{enumerate}
\item \emph{Small coupling at the operator level.} \eqref{eq:kstatic} is used
  with the uncoupled denominator only where $k^{2}\ll1$; here
  $k^{2}=0.5581$, which is \emph{not} small, so the coupled denominator is
  used and the uncoupled form appears only in the comparison with the
  closed form, where it is stated.
\item \emph{Small window.} The truncation \eqref{eq:Mdef} is written for a
  window small against the length over which the strain gradient varies, and
  for the strip of Table~\ref{tab:params} we take $a/L\le0.2$. That is an
  assumed operating range, not a validated truncation threshold: with no
  experiment and no exact averaging operator available, nothing here fixes
  where the truncation fails. What Fig.~\ref{fig:window}(a) shows is that
  inside the range the retained second moment reproduces the exact average to
  a fraction of a per cent on every mesh, which makes the model
  \emph{self-consistent} there and is all that is claimed.
\item \emph{Operating band.} The damping map is claimed for $\omega\tau\le1$,
  the band in which the spring-pot relaxes around the device frequency. Beyond
  it the map extrapolates the fitted law, and Table~\ref{tab:modulus} shows how
  far the integer-order alternative departs.
\end{enumerate}

\FloatBarrier
\subsection{Limitations}
\label{sec:limitations}

\emph{Two dimensions and unverified inputs.} The formulation is
two-dimensional with a plane-stress reduction, which is a modeling choice: plane
strain would change the extensional coupling, and the three-dimensional
solution differs from both, so the thin-strip results of
Sec.~\ref{sec:thickness} are two-dimensional results consistent with
Table~\ref{tab:params}. Two inputs are nominal. The density
$\rho=\SI{6020}{\kilo\gram\per\cubic\metre}$ rescales every frequency by a
common factor, so the kilohertz numbers of Secs.~\ref{sec:orientation}
and~\ref{sec:application} would move while the ratios stand; and
$\mu=10^{-6}$\,\si{\coulomb\per\metre} scales every flexoelectric coupling as
$\mu^{2}$, so the growth factor, the saturation and the comparison with the
piezoelectric response are unaffected while the absolute coupling is not.

\emph{What the verification does not establish.} The two-significant-digit
statement of Sec.~\ref{sec:thickness} is a limitation as much as a result: the
third digit there is not a property of the model, and the same holds, less
visibly, for any difference of nearly equal eigenvalues. The verification of
Secs.~\ref{sec:verification} and~\ref{app:verification} shows that the
implementation solves the stated equations; it does not show that they describe
a particular material, which needs the measurements of
Secs.~\ref{sec:calib_a_mu} and~\ref{sec:calib_alpha_tau}.

\FloatBarrier
